\documentclass[french]{book}
\usepackage{teach}
\usepackage{amsmath,amsfonts,amssymb}
\usepackage{pgfplots}
%\usepackage{geometry}
%\geometry{top=1cm, bottom=1cm, left=2cm, right=2cm}
\usepackage[utf8]{inputenc}

%\usepackage[latin1]{inputenc}
%\usepackage[T1]{fontenc}

\usepackage[french]{babel}
\redefinecolor{thm}{title}{bgcolor}{alizarin}
\redefinecolor{prop}{title}{bgcolor}{meatbrown}
\redefinecolor{defn}{title}{bgcolor}{olivine}
\definecolor{alizarin}{rgb}{0.82, 0.1, 0.26}
\definecolor{meatbrown}{rgb}{0.9, 0.72, 0.23}
\definecolor{olivine}{rgb}{0.6, 0.73, 0.45}
\usepackage{pgf,tikz}

\usepackage{mathrsfs}
\usetikzlibrary{arrows}
\usetikzlibrary{patterns}

\begin{document}
\thispagestyle{empty}

\section*{Devoir surveillé III}

Voici les formules du cours: 
\begin{itemize}
\item Proportion d'une sous population qui contient $n$ individus  par rapport à une population de $N$ individus $\frac{n}{N}$.
\item $C_m=1+t$
\item $V_d \times C_m= V_a $
\item $C_{mg}=C_{m_1} \times C_{m_2} \times \cdots \times C_{m_p}$
\item $C_{mg}=1+T$
\item $C_{m'}=\frac{1}{C_m}$
\item $V_a \times C_{m'}= V_d $\\

\end{itemize}

\underline{Exercice 1}\\


Dans un lycée de 500 élève, une classe représente 7\% des élèves du lycée. Dans cette classe, 14 élèves étudient uniquement l’anglais, 12 élèves étudient uniquement l’espagnol et 5 élèves étudient les deux langue.\\


1) Quel est la proportion d'élève qui étudie l'espagnol par rapport à la classe?\\

2) Quel est la proportion d'élève qui étudie l'anglais parmis ceux qui étudie aussi l'espagnol?\\

3) Quel est la proportion d'élève qui étudie l'anglais et l'espagnol par rapport à la classe? Donner deux méthodes, une à partir de l'énoncé et une à partir des question 1) et 2).\\

4) Quel est la proportion d'élève qui n'étudie aucune des deux langues ?\\

\underline{Exercice 2}\\

Du plâtre pour faire des moulures au plafond est vendu dans un magasin. Kévin, un vendeur du magasin, applique une réduction de 20\% sur le plâtre. Il décide 2 semaines plus tard d'appliquer une augmentation de 20\% sur le plâtre pour remettre le plâtre au prix initial. \\

5) Expliquer à l'aide d'un calcul, si Kévin à raison ou non.\\

Inès, une collègue de Kévin, lui dit que les calculs ne sont pas bon. Par suite elle dit qu'il faut multiplier le prix réduit par 1,25 pour obtenir le prix initial. \\

6) Justifier par un calcul, si Inès a raison ou non.\\

\underline{Exercice 3} \\ 

7) Déterminer le coefficient multiplicateur global d'une réduction de 20\% puis d'une augmentation de 10\% d'une réduction de 25\% et d'une augmentation de 50\%. \\

8) En déduire le taux d'évolution global. Ainsi cela correspond à une ................... de ..... \%.\\

\underline{Exercice 4} \\

Un téléphone à subit 30\% de réduction, son prix après réduction est de 811,3 euros.

9) Quel était le prix de départ du téléphone ? \\

Un ordinateur à subit 30\% de réduction puis une augmentation de 20\%, son prix après ces application de réducion et d'augmentation est de 504 euros.\\

10) Quel était le prix de départ de l'ordinateur?
 
 

\end{document}